\section*{Capítulo \ref{ch:g10:lines} --- Equações de retas e sistemas lineares}

\begin{solution}{exo:g10:lines:1}
$y = 2x - 3$: \omterm{def:g10:reffunc:affine}{coeficiente angular} $2$, \omterm{def:g10:reffunc:affine}{coeficiente linear} $-3$.
$y = -\frac13 x + 1$: \omterm{def:g10:reffunc:affine}{coeficiente angular} $-\frac13$, \omterm{def:g10:reffunc:affine}{coeficiente linear}
$1$.
$y = 4$: \omterm{def:g10:reffunc:affine}{coeficiente angular} $0$, \omterm{def:g10:reffunc:affine}{coeficiente linear} $4$ (reta
horizontal).
$x = -2$: reta vertical, sem \omterm{def:g10:reffunc:affine}{coeficiente angular} e sem \omterm{thm:g10:lines:reduced}{equação reduzida}.
\end{solution}

\begin{solution}{exo:g10:lines:2}
$3 \times 2 + 1 = 7$: sim, $A$ está na reta.
$3 \times (-1) + 1 = -2 \neq -1$: $B$ não está.
$3 \times 4 + 1 = 13$: sim, $C$ está na reta.
\end{solution}

\begin{solution}{exo:g10:lines:3}
(a) $m = \frac{8 - 2}{3 - 0} = 2$ e $p = y_A = 2$ (o ponto $A$ está sobre
o eixo $y$): $y = 2x + 2$.

(b) $m = \frac{-4 - 5}{2 - (-1)} = \frac{-9}{3} = -3$; então
$5 = -3 \times (-1) + p$ dá $p = 2$: $y = -3x + 2$.

(c) $x_A = x_B = 4$: reta vertical $x = 4$.
\end{solution}

\begin{solution}{exo:g10:lines:4}
$y = \frac{6x+2}{2} = 3x + 1$. As retas de \omterm{def:g10:reffunc:affine}{coeficiente angular} $3$ são
$y = 3x - 1$, $y = 3x + 5$ e $y = 3x + 1$: essas três são paralelas entre
si; nenhuma delas é igual a outra (\omterm{def:g10:reffunc:affine}{coeficientes lineares} diferentes), e
nenhuma é paralela a $y = -3x + 2$.
\end{solution}

\begin{solution}{exo:g10:lines:5}
Primeiro sistema: substitua $y = 2x - 1$ em $3x + y = 9$:
$3x + 2x - 1 = 9$, logo $5x = 10$, $x = 2$, e então $y = 3$. Solução:
$(2, 3)$.

Segundo sistema: de $x - y = 4$, $x = y + 4$. Então
$2(y + 4) + 3y = 3$, logo $5y = -5$, $y = -1$, e então $x = 3$. Solução:
$(3, -1)$.
\end{solution}

\begin{solution}{exo:g10:lines:6}
Primeiro sistema: somar as \omterm{def:g10:algebra:equation}{equações} elimina $y$: $2x = 14$, $x = 7$; e
então $7 + y = 10$ dá $y = 3$. Solução: $(7, 3)$.

Segundo sistema: multiplique a primeira \omterm{def:g10:algebra:equation}{equação} por $5$ e a segunda por
$-2$:
\[
\begin{cases}
15x + 10y = 60 \\
-4x - 10y = -16
\end{cases}
\]
Somando: $11x = 44$, logo $x = 4$; e então $3 \times 4 + 2y = 12$ dá
$y = 0$. Solução: $(4, 0)$.
\end{solution}

\begin{solution}{exo:g10:lines:7}
Calcule $ab' - a'b$ em cada caso.

Primeiro: $2 \times 2 - (-4) \times (-1) = 4 - 4 = 0$, e a segunda
\omterm{def:g10:algebra:equation}{equação} é $-2$ vezes a primeira: duas vezes a mesma reta, uma infinidade
de soluções.

Segundo: $ab' - a'b = 0$ de novo, mas os lados direitos não estão na
razão $-2$ ($1 \neq -6$): duas retas paralelas distintas, nenhuma
solução.

Terceiro: $2 \times 1 - 1 \times (-1) = 3 \neq 0$: exatamente uma
solução.
\end{solution}

\begin{solution}{exo:g10:lines:8}
Seja $x$ o número de ingressos infantis e $y$ o de ingressos de adulto:
\[
\begin{cases}
x + y = 200 \\
5x + 9y = 1352 .
\end{cases}
\]
Da primeira \omterm{def:g10:algebra:equation}{equação}, $x = 200 - y$; substituindo,
$5(200 - y) + 9y = 1352$, logo $1000 + 4y = 1352$, $y = 88$, e então
$x = 112$. Portanto, $112$ ingressos infantis e $88$ de adulto.
Verificação: $5 \times 112 + 9 \times 88 = 560 + 792 = 1352$.
\end{solution}

\begin{solution}{exo:g10:lines:9}
\emph{1.} Paralela significa mesmo \omterm{def:g10:reffunc:affine}{coeficiente angular} $2$:
$y = 2x + p$ com $4 = 2 \times 1 + p$, logo $p = 2$:
$d'\colon y = 2x + 2$.

\emph{2.} Sobre o eixo $x$, $y = 0$: $2x + 2 = 0$ dá $x = -1$. O ponto
de \omterm{def:g10:numbers:interunion}{interseção} é $(-1, 0)$.
\end{solution}

\begin{solution}{exo:g10:lines:10}
\emph{1.} \omterm{prop:g10:coordgeom:midpoint}{Ponto médio} de $[BC]$: $A' = \left(\frac{6+2}{2},
\frac{2+4}{2}\right) = (4, 3)$. A mediana que parte de $A(0,0)$ tem
\omterm{def:g10:reffunc:affine}{coeficiente angular} $\frac{3 - 0}{4 - 0} = \frac34$: \omterm{def:g10:algebra:equation}{equação}
$y = \frac34 x$.

\omterm{prop:g10:coordgeom:midpoint}{Ponto médio} de $[AC]$: $B' = (1, 2)$. A mediana que parte de $B(6, 2)$
tem \omterm{def:g10:reffunc:affine}{coeficiente angular} $\frac{2 - 2}{1 - 6} = 0$: é a reta horizontal
$y = 2$.

\emph{2.} \omterm{def:g10:numbers:interunion}{Interseção}: $\frac34 x = 2$ dá $x = \frac83$, logo
$G\left(\frac83, 2\right)$. A terceira mediana liga $C(2,4)$ ao \omterm{prop:g10:coordgeom:midpoint}{ponto
médio} de $[AB]$, $C' = (3, 1)$; seu \omterm{def:g10:reffunc:affine}{coeficiente angular} é
$\frac{1 - 4}{3 - 2} = -3$, \omterm{def:g10:algebra:equation}{equação} $y = -3x + 10$. Em $x = \frac83$:
$y = -8 + 10 = 2$: $G$ está sobre ela. E, de fato,
$G = \left(\frac{0 + 6 + 2}{3}, \frac{0 + 2 + 4}{3}\right)$: o
baricentro é a média das \omterm{def:g10:coordgeom:system}{coordenadas} dos vértices.
\end{solution}

\begin{solution}{pb:g10:lines:1}
\textbf{1.} \omterm{def:g10:reffunc:affine}{Coeficiente angular} $\frac{8 - 2}{3 - 1} = 3$;
passando por $(1, 2)$: $y = 3x - 1$.

\textbf{2.} $y = 3x - 1$ e $y = 3x + 4$: mesmo \omterm{def:g10:reffunc:affine}{coeficiente
angular}, paralelas (e distintas). Cruzando $y = 3x - 1$ com
$y = -x + 7$: $3x - 1 = -x + 7$ dá $x = 2$, $y = 5$: ponto
$(2, 5)$.

\textbf{3.} Os dois pontos têm $x = 2$: a reta é vertical, de
\omterm{def:g10:algebra:equation}{equação} $x = 2$. Uma forma reduzida $y = mx + p$ responde ``um
$y$ para cada $x$'', o que uma reta vertical recusa: seu
\omterm{def:g10:reffunc:affine}{coeficiente angular} seria infinito.

\textbf{4.} $3 \times 4 - 1 = 11$: sim, $(4, 11)$ está na reta.
$3 \times (-1) - 1 = -4 \neq -5$: não.

\textbf{5.} $y = -2x + 3$; corta os eixos em $(0, 3)$ e
$\left(\frac32, 0\right)$. Área do triângulo:
$\frac12 \times \frac32 \times 3 = \frac94 = 2.25$.

\textbf{6.} Substituindo: $3x + 2x - 3 = 7$, $x = 2$,
$y = 1$: solução $(2, 1)$.

\textbf{7.} Somando: $7x = 7$, $x = 1$; então $3y = 6$, $y = 2$:
solução $(1, 2)$.

\textbf{8.} Primeiro sistema: a segunda \omterm{def:g10:algebra:equation}{equação} é o dobro da
primeira: uma reta contada duas vezes --- uma infinidade de
soluções (todos os pontos de $2x - y = 3$). Segundo sistema:
lados esquerdos proporcionais, lados direitos não
($2 \times 3 = 6 \neq 1$): duas retas paralelas --- nenhuma
solução. Tricotomia: \omterm{def:g10:reffunc:affine}{coeficientes angulares} distintos $\to$ um
cruzamento; \omterm{def:g10:reffunc:affine}{coeficientes angulares} iguais e \omterm{def:g10:reffunc:affine}{coeficientes
lineares} diferentes $\to$ paralelas, nenhuma solução; tudo
igual $\to$ a mesma reta, uma infinidade.

\textbf{9.} Os \omterm{def:g10:vectors:vector}{vetores} diretores das retas são $(-b, a)$ e
$(-d, c)$ (ou, equivalentemente, os \omterm{def:g10:reffunc:affine}{coeficientes angulares}
$-\frac ab$ e $-\frac cd$), \omterm{def:g10:vectors:collinear}{colineares} exatamente quando
$ad - bc = 0$ --- o determinante do \cref{pb:g10:vectors:1} em
novo emprego. Valores: questão 7:
$2 \times (-3) - 3 \times 5 = -21 \neq 0$: solução única.
Questão 8: $2 \times (-2) - (-1) \times 4 = 0$ nos dois casos:
paralelas ou idênticas, e as constantes separam os dois casos.

\textbf{10.} $ad - bc = 6 - 2m$: nulo para $m = 3$. Aí
$x + 3y = 3$ contra $2x + 6y = 7$: dobrar a primeira dá
$2x + 6y = 6 \neq 7$: paralelas, \emph{nenhuma solução}. Para
$m \neq 3$, a eliminação dá a solução única
$y = \frac{1}{6 - 2m}$, e então $x = 3 - \frac{m}{6 - 2m}$.

\textbf{11.} $f + c = 35$, $2f + 4c = 94$: dividindo por dois,
$f + 2c = 47$, logo $c = 12$ e $f = 23$: vinte e três faisões e
doze coelhos --- a própria resposta dos \emph{Nove Capítulos}.

\textbf{12.} $x + y = 30$ e $0.6x + 0.2y = 0.35 \times 30 =
10.5$: subtraindo $0.2 \times$ a primeira, $0.4x = 4.5$:
$x = 11.25$ L de xarope e $y = 18.75$ L de bebida.

\textbf{13.} Algarismos $a$, $b$: $a + b = 12$ e
$(10a + b) - (10b + a) = 9(a - b) = 18$, logo $a - b = 2$:
$a = 7$, $b = 5$: o número é $75$.

\textbf{14.} Somando: $5c + 5r = 18$, logo $c + r = 3.60$ ---
um café mais um croissant. Subtraindo: $c - r = 0.20$. Daí
$c = 1.90$ e $r = 1.70$ euros: o atalho simétrico resolveu tudo
sem nenhuma substituição.

\textbf{15.} Dobrar o primeiro anúncio dá $4$ maçãs $+ 2$
peras $= 10$ euros, mas a banca cobra $9$: o sistema é
incompatível --- retas paralelas, conjunto solução vazio.
Veredicto: os dois preços anunciados não podem estar ambos
certos; ou há um desconto escondido, ou há um erro de conta.

\textbf{16.} Na origem: $0 > 3 \times 0 - 1 = -1$: verdadeiro,
logo $(0,0)$ está na região $y > 3x - 1$ (acima da reta). Teste
qualquer ponto substituindo-o: acima ou abaixo conforme a
desigualdade seja verdadeira ou falsa.

\textbf{17.} Vértices: $(0, 0)$; $(4, 0)$ (eixos e
$x + y = 4$); $(0, 1)$ (eixo e $y = 2x + 1$); e $x + y = 4$ com
$y = 2x + 1$: $x = 1$, $y = 3$: $(1, 3)$.

\textbf{18.} $P(0,0) = 0$; $P(4, 0) = 12$; $P(0, 1) = 2$;
$P(1, 3) = 9$. \omterm{def:g10:functions:extrema}{Máximo} no vértice $(4, 0)$: o plano $x = 4$,
$y = 0$ rende $12$ --- os vértices mandam, como as retas
paralelas $3x + 2y = c$ deixam ver ao varrer o plano.

\textbf{19.} $1.5x + 3 = 2x$ dá $x = 6$, $y = 12$: em seis
quilômetros os dois táxis cobram doze euros --- o cruzamento de
equilíbrio do problema de fim de semana sobre os dois
termômetros, no volume do ensino fundamental, agora como
solução única de um sistema.

\textbf{20.} As soluções de uma \omterm{def:g10:algebra:equation}{equação} linear desenham uma
reta; um sistema superpõe duas retas, e $ad - bc$ diz de
imediato se elas se cruzam uma vez (não nulo) ou caem nos casos
paralelas-ou-idênticas (nulo). Empilhe inequações em vez disso
e as soluções formam uma região poligonal cujos vértices
carregam qualquer ótimo linear. A álgebra linear generaliza o
determinante para um número qualquer de incógnitas; a
programação linear industrializa os vértices.
\end{solution}
